\chapter{Theoretical Background}

This chapter presents the theoretical foundations for constructing Zero-Knowledge Proof (ZKP) systems for verifying the  logistic regression predictions and Accuracy. We outline the core concepts including zero-knowledge proofs, arithmetic circuits, polynomial commitment schemes, sigmoid approximation techniques, and logistic regression model formulation.

\section{Zero-Knowledge Proofs}

\begin{definition}\label{def:zkp}
\textbf{Zero-Knowledge Proof (ZKP):} A ZKP is an interactive or non-interactive cryptographic protocol enabling a prover to convince a verifier that a statement is true, without revealing any additional information beyond the truth of the statement itself~\cite{goldwasser1989knowledge}.
\end{definition}

A ZKP system satisfies three essential properties:  
\textbf{Completeness}, where an honest prover convinces the verifier;  
\textbf{Soundness}, where a cheating prover cannot convince the verifier of a false statement; and  
\textbf{Zero-Knowledge}, where the proof reveals nothing beyond the validity of the statement.

\begin{remark}
Modern proof systems such as SNARKs and PLONK extend these principles to arbitrary computations by encoding them as arithmetic circuits~\cite{kates2020plonk}.
\end{remark}

\section{Arithmetic Circuits and Constraint Systems}

\begin{definition}\label{def:circuits}
\textbf{Arithmetic Circuit:} An arithmetic circuit represents a computation using addition and multiplication gates over a finite field, enabling structured modeling of machine learning computations within zero-knowledge systems~\cite{gennaro2013quadratic}.
\end{definition}

For a computation $f(x, w) = y$, the circuit expresses:
\begin{equation}
\text{Constraints}(x, w) = 0,
\end{equation}
where $x$ denotes public inputs and $w$ represents private witness values.

\begin{remark}
Arithmetic circuits form the foundation for expressing logistic regression operations, including linear transformations and sigmoid approximations, within ZK systems.
\end{remark}

\section{Polynomial Commitment Schemes}

\begin{definition}\label{def:polycommit}
\textbf{Polynomial Commitment Scheme:} A cryptographic primitive that allows a prover to commit to a polynomial and later prove evaluations at specific points without revealing the polynomial itself~\cite{kates2020plonk}.
\end{definition}

The commitment is represented as:
\begin{equation}
    C = \text{Commit}(f(x)),
\end{equation}
and the prover later provides
\[
\text{Proof}(k, f(k)),
\]
allowing the verifier to confirm correctness without accessing the entire polynomial.

\begin{remark}
Polynomial commitments enable succinct verification and efficient batching of constraints, essential for scalability in proof systems.
\end{remark}

\section{Sigmoid Function and Approximation}

\begin{definition}\label{def:sigmoid}
\textbf{Sigmoid Function:} The sigmoid function in logistic regression is defined as:
\begin{equation}
    \sigma(z) = \frac{1}{1 + e^{-z}}.
\end{equation}
\end{definition}

Direct computation of $\sigma(z)$ is expensive within ZKP systems because of the exponential and division operations. Hence, approximation methods are employed~\cite{cryptoeprint:2025/507}.

Common ZKP-friendly approximations include:
\begin{itemize}
    \item Piecewise linear approximation,
    \item Polynomial approximation (e.g., Taylor or Chebyshev series),
    \item Lookup-table-based approximation.
\end{itemize}

\begin{remark}
Lookup-table-based approximations significantly reduce circuit complexity and proving time while maintaining adequate numerical precision.
\end{remark}

\section{Logistic Regression Model}

\begin{definition}\label{def:logreg}
\textbf{Logistic Regression:} A statistical model for binary classification that maps input features to probabilities using:
\begin{equation}
    z = w^T x + b,
\end{equation}
where $w$ is the weight vector, $x$ is the input vector, and $b$ is the bias term~\cite{hosmer2013applied}.
\end{definition}

The predicted probability is calculated using the sigmoid function:
\begin{equation}\label{eq:prediction}
\widehat{y} = \sigma(z) = \frac{1}{1 + e^{-z}}.
\end{equation}

\begin{remark}
Within a ZKP framework, computing $z$ is straightforward using arithmetic circuits, while evaluating $\sigma(z)$ requires approximation optimizations to balance efficiency and accuracy~\cite{cryptoeprint:2023/1174}.
\end{remark}

\clearpage
